\documentclass[12pt,a4paper]{article}

% ── Paquetes ──────────────────────────────────────────────────────────────────
\usepackage[utf8]{inputenc}
\usepackage[T1]{fontenc}
\usepackage[spanish]{babel}
\usepackage[margin=2.5cm]{geometry}
\usepackage{booktabs}
\usepackage[table]{xcolor}
\usepackage{fancyhdr}
\usepackage{titlesec}
\usepackage[many, breakable]{tcolorbox}
\usepackage{graphicx}
\usepackage{amssymb}
\usepackage{float}
\usepackage{enumitem}
\usepackage{hyperref}

\tcbset{
  cajainfo/.style={enhanced,breakable,colback=verdeclaro,colframe=verdeosc,
    fonttitle=\bfseries\small\color{white},coltitle=white,
    attach boxed title to top left={yshift=-2mm,xshift=4mm},
    boxed title style={colback=verdeosc},left=6pt,right=6pt,top=10pt,bottom=6pt}
}

\definecolor{verdeosc}{RGB}{26,107,60}
\definecolor{verdemid}{RGB}{46,139,87}
\definecolor{dorado}{RGB}{212,172,13}
\definecolor{azulclaro}{RGB}{236,240,247}
\definecolor{verdeclaro}{RGB}{232,244,232}
\definecolor{verde}{RGB}{40,167,69}
\definecolor{rojo}{RGB}{220,53,69}

\setlength{\headheight}{30pt}
\pagestyle{fancy}
\fancyhf{}
\fancyhead[L]{\small\color{verdemid}SICVEC 2026}
\fancyhead[R]{\small\color{verdemid}Manual de Memorias}
\fancyfoot[C]{\small\thepage}
\renewcommand{\headrulewidth}{0.6pt}
\renewcommand{\footrulewidth}{0.4pt}

\titleformat{\section}
  {\normalfont\large\bfseries\color{white}}
  {\fcolorbox{verdeosc}{verdeosc}{\thesection}}{0.7em}{}
  [\vspace{-2pt}\color{verdeosc}\rule{\linewidth}{0.8pt}]

\titleformat{\subsection}
  {\normalfont\normalsize\bfseries\color{verdemid}}{\thesubsection}{0.6em}{}

% ============================================================================
\begin{document}
% ============================================================================

\thispagestyle{empty}

\begin{center}
{\Large\bfseries\color{verdeosc} SIMPOSIO INTERNACIONAL DE CIENCIA VERDE Y ECONOMÍA CIRCULAR}\\
{\normalsize\color{verdemid} SICVEC 2026 -- 19 y 20 de octubre}\\[1.5cm]
{\LARGE\bfseries MANUAL DE MEMORIAS Y PROCESO}\\
{\large\color{verdemid} Recepción de Resúmenes -- Peer Review -- Publicación}\\
\end{center}

\vspace{2cm}

\begin{tcolorbox}[enhanced, colback=verdeclaro, colframe=verdeosc, boxrule=2pt, left=15pt, right=15pt, top=15pt, bottom=15pt]

\noindent\textbf{\Large Objetivo del Manual}

Este documento orienta a investigadores, estudiantes y profesionales sobre cómo participar en SICVEC 2026 con sus trabajos científicos, presentando resúmenes que serán sometidos a evaluación de pares (peer review) con estándares internacionales de calidad.

\vspace{0.5cm}

\noindent\textbf{Documentos relacionados:}
\begin{itemize}[label={\color{verdemid}\checkmark}]
  \item Guía de Presentación de Resúmenes
  \item Template de Resumen Científico (para descargar y completar)
  \item Rúbrica de Evaluación (usado por revisores)
\end{itemize}

\end{tcolorbox}

\newpage
\tableofcontents
\newpage

% ============================================================================
\section{Visión General del Proceso}
% ============================================================================

El proceso de gestión de memorias de SICVEC 2026 sigue estándares internacionales de publicaciones científicas:

\begin{tcolorbox}[enhanced, colback=white, colframe=verdeosc, boxrule=1.5pt]

\textbf{Flujo del Proceso:}

\begin{enumerate}[label={\textcolor{verdemid}{\textbf{Paso \arabic{*}}}}]
  \item \textbf{Convocatoria Abierta} (22 de septiembre) -- Investigadores son invitados a participar
  \item \textbf{Envío de Resúmenes} (hasta 4 de octubre) -- Autores cargan sus resúmenes
  \item \textbf{Admisibilidad y Asignación de Revisores} (4-5 de octubre) -- Verificación de formato y 2 especialistas por resumen
  \item \textbf{Evaluación de Pares} (5-7 de octubre) -- Revisores califican en línea
  \item \textbf{Decisión Académica} (8 de octubre) -- Comité científico decide
  \item \textbf{Notificación} (8-9 de octubre) -- Email con resultado y feedback
  \item \textbf{Material Final} (hasta 14 octubre) -- Autores aceptados envían versión definitiva
  \item \textbf{Publicación en Memorias} (20 octubre) -- Disponible durante el evento
  \item \textbf{Archivamiento Digital} (post-evento) -- Acceso permanente en plataforma
\end{enumerate}

\end{tcolorbox}

% ============================================================================
\section{Cronograma Detallado}
% ============================================================================

\begin{table}[h]
\centering
\small
\begin{tabular}{|l|l|l|}
\hline
\rowcolor{verdeclaro}
\textbf{Fecha} & \textbf{Actividad} & \textbf{Responsable} \\
\hline
22 de septiembre & Apertura de convocatoria & María Ximena Díaz \\
\hline
22 sept - 4 oct & Período de envío abierto & Autores \\
\hline
4 de octubre 23:59 & CIERRE de recepción & Sistema automático \\
\hline
4-5 de octubre & Revisión admisibilidad y asignación de revisores & Coord. Académica \\
\hline
5-7 de octubre & Evaluación de pares & Revisores (2 por trabajo) \\
\hline
8 de octubre & Análisis de evaluaciones y decisiones finales & Comité Científico / Coord. Académica \\
\hline
8-9 de octubre & Notificación a autores & Sistema de email \\
\hline
hasta 14 de octubre & Envío de material final & Autores aceptados \\
\hline
15 de octubre & Límite de pago de inscripciones & Autores y participantes \\
\hline
19-20 octubre & SIMPOSIO & Todos \\
\hline
20 octubre & Publicación de memorias & Coord. General \\
\hline
\end{tabular}
\end{table}

% ============================================================================
\section{Tipos de Trabajos Aceptados}
% ============================================================================

\subsection{1. Ponencia Oral (20 minutos)}

Presentación de investigaciones completas con resultados concluyentes.

\textbf{Resumen:} Máx. 300 palabras | \textbf{Público:} Investigadores avanzados, profesionales

\subsection{2. Poster Científico (formato: 90×120 cm)}

Proyectos en desarrollo, resultados preliminares, experiencias innovadoras.

\textbf{Resumen:} Máx. 250 palabras | \textbf{Público:} Estudiantes, investigadores junior

\subsection{3. Taller Práctico (90 minutos)}

Capacitación hands-on, case studies, demostraciones tecnológicas.

\textbf{Propuesta:} Máx. 500 palabras | \textbf{Capacidad:} 25-40 participantes

% ============================================================================
\section{Criterios de Evaluación Científica}
% ============================================================================

Todos los resúmenes son evaluados mediante rúbrica de 5 criterios con ponderación:

\begin{table}[H]
\centering
\begin{tabular}{|l|r|l|}
\hline
\rowcolor{verdeclaro}
\textbf{Criterio} & \textbf{Peso} & \textbf{¿Qué se evalúa?} \\
\hline
Relevancia & 25\% & Pertinencia a SICVEC, aporte a ciencia verde \\
\hline
Originalidad & 25\% & Ideas nuevas, innovación, estado del arte \\
\hline
Calidad Científica & 25\% & Rigor metodológico, validez de resultados \\
\hline
Claridad & 15\% & Redacción clara, estructura lógica \\
\hline
Impacto Potencial & 10\% & Aplicabilidad, contribución a sostenibilidad \\
\hline
\end{tabular}
\end{table}

\textbf{Escala de calificación:} 1 (Insuficiente) a 5 (Excelente)

\textbf{Puntuación mínima para aceptación:} 3.5/5.0

% ============================================================================
\section{Sistema de Peer Review}
% ============================================================================

\subsection{Características del Proceso}

\begin{itemize}[label={\textcolor{verdemid}{\checkmark}}]
  \item \textbf{Double-Blind:} Identidades de autores y revisores confidenciales
  \item \textbf{Especializados:} Revisores son expertos en el tema
  \item \textbf{Transparente:} Feedback constructivo para mejorar trabajos
  \item \textbf{Justo:} Criterios iguales para todos los trabajos
  \item \textbf{Oportuno:} Decisiones en menos de 2 meses
\end{itemize}

\subsection{Rol de los Revisores}

Los revisores son investigadores y profesionales calificados que:

\begin{enumerate}
  \item Evalúan objetivamente usando rúbrica estandarizada
  \item Proporcionan feedback constructivo (comentarios detallados)
  \item Recomiendan aceptación, aceptación con cambios, o rechazo
  \item Respetan confidencialidad (revisión ciega)
  \item Declaran ausencia de conflicto de interés
\end{enumerate}

\subsection{Posibles Decisiones}

\begin{tcolorbox}[enhanced, colback=green!10, colframe=verde, boxrule=1pt, left=10pt, right=10pt, top=8pt, bottom=8pt]
\textbf{ACEPTADO} -- Trabajo cumple con estándares de calidad. Será incluido en programa oficial del simposio y memorias.
\end{tcolorbox}

\begin{tcolorbox}[enhanced, colback=yellow!15, colframe=dorado, boxrule=1pt, left=10pt, right=10pt, top=8pt, bottom=8pt]
\textbf{ACEPTADO CON CAMBIOS} -- Se solicitan cambios menores. Autores deben corregir antes de 10 nov. para inclusion en memorias.
\end{tcolorbox}

\begin{tcolorbox}[enhanced, colback=red!10, colframe=rojo, boxrule=1pt, left=10pt, right=10pt, top=8pt, bottom=8pt]
\textbf{RECHAZADO} -- Trabajo no cumple criterios de calidad/relevancia. Se proporciona feedback para futuras mejoras.
\end{tcolorbox}

% ============================================================================
\section{Memorias Digitales}
% ============================================================================

Las memorias de SICVEC 2026 incluirán:

\subsection{Contenidos}

\begin{itemize}[label={\textcolor{verdemid}{\textbullet}}]
  \item Resúmenes de todas las ponencias aceptadas (formato estandarizado)
  \item Posters científicos en PDF (de alta calidad)
  \item Biografías de conferencistas internacionales
  \item Programa científico detallado (cronograma, horarios, salas)
  \item Listado completo de participantes
  \item Índices temáticos (por eje, por institución, por autor)
  \item Documentos de política pública generados
\end{itemize}

\subsection{Formato y Distribución}

\begin{itemize}[label={\textcolor{verdemid}{\checkmark}}]
  \item \textbf{Formato primario:} PDF interactivo descargable
  \item \textbf{Acceso:} Plataforma en línea de SICVEC (acceso permanente)
  \item \textbf{ISBN:} Memorias tendrán número de depósito legal
  \item \textbf{Citas:} Se puede citar como: ``Memorias de SICVEC 2026''
  \item \textbf{Idioma:} Inglés y español (abstracts bilingües preferiblemente)
\end{itemize}

\subsection{Derechos de Autor}

\begin{tcolorbox}[cajainfo, title=Propiedad Intelectual]
\begin{itemize}
  \item Los \textbf{autores retienen} todos los derechos sobre sus trabajos
  \item SICVEC 2026 tiene derecho a reproducir resúmenes en memorias
  \item Se requiere \textbf{consentimiento escrito} para usos posteriores
  \item Los resúmenes \textbf{no pueden duplicarse} en otros eventos simultáneos
  \item Citación apropiada debe indicar ``Memorias SICVEC 2026''
\end{itemize}
\end{tcolorbox}

% ============================================================================
\section{Oportunidades Post-Evento}
% ============================================================================

\subsection{Publicación en Revista Indexada}

Se está coordinando con editores para publicar artículos completos en:

\begin{itemize}[label={\textcolor{verdemid}{\checkmark}}]
  \item Revista indexada en Scopus/WoS de sostenibilidad
  \item Fast-track para autores de SICVEC 2026
  \item Descuento en cuotas de publicación
\end{itemize}

\subsection{Libro de Memorias Arbitrado}

\begin{itemize}[label={\textcolor{verdemid}{\checkmark}}]
  \item Selección de 15-20 mejores resúmenes
  \item Expansión a artículos completos (5,000-8,000 palabras)
  \item ISBN y depósito legal
  \item Distribución nacional e internacional
\end{itemize}

\subsection{Repositorio Digital Permanente}

\begin{itemize}[label={\textcolor{verdemid}{\checkmark}}]
  \item Acceso abierto a todos los resúmenes aceptados
  \item Búsqueda por tema, autor, institución
  \item DOI para cada trabajo (citabilidad)
  \item Integración con Google Scholar y ResearchGate
\end{itemize}

% ============================================================================
\section{Preguntas Frecuentes}
% ============================================================================

\begin{tcolorbox}[enhanced, colback=azulclaro, colframe=verdemid, boxrule=1.5pt, left=10pt, right=10pt, top=10pt, bottom=10pt]

\noindent\textbf{P: ¿Cuál es el límite de palabras para un resumen?}\\
\textbf{R:} 300 palabras para ponencias, 250 para posters (excluye referencias).

\vspace{0.5cm}

\noindent\textbf{P: ¿Puedo enviar mi resumen en otro idioma además de español?}\\
\textbf{R:} Sí, aceptamos español e inglés. Preferiblemente con resumen bilingüe.

\vspace{0.5cm}

\noindent\textbf{P: ¿Quién puede revisar mi trabajo?}\\
\textbf{R:} Especialistas calificados del comité científico (autores desconocen identidad).

\vspace{0.5cm}

\noindent\textbf{P: ¿Recibiré feedback si rechazan mi trabajo?}\\
\textbf{R:} Sí, los revisores proporcionan comentarios constructivos para mejora.

\vspace{0.5cm}

\noindent\textbf{P: ¿Mis datos estarán confidenciales durante peer review?}\\
\textbf{R:} Totalmente. Sistema double-blind protege privacidad de autores y revisores.

\end{tcolorbox}

% ============================================================================
\section{Contacto y Soporte}
% ============================================================================

Para preguntas sobre el proceso de memorias y peer review:

\vspace{0.8cm}

\begin{center}
\begin{tabular}{ll}
\textbf{Coordinadora Académica:} & María Ximena Díaz \\
\textbf{Email:} & sicvec.academico@unisucre.edu.co \\
\textbf{Teléfono:} & [Completar] \\
\textbf{Horario:} & Lunes-viernes, 8:00-17:00 (Hora Colombia) \\
\end{tabular}
\end{center}

\vspace{1cm}

\textbf{Plataforma de inscripción:} [URL - Completar]

\textbf{Sitio web del evento:} www.sicvec2026.unisucre.edu.co

\end{document}
